% =============================================================================
%  第 1 章  线性方程组与消元法
% =============================================================================
\chapter{线性方程组与消元法}\label{chap:systems}

\epigraph{以右行上禾遍乘中行，而以直除。}{《九章算术 · 方程》}
\chaptertoc

求解线性方程组是线性代数最基本的问题。本章介绍解线性方程组的消元法，引入矩阵和初等行
变换的语言，并完整地回答两个问题：方程组什么时候有解？有解时，解是否唯一？

\section{从《九章算术》说起}\label{sec:jiuzhang}

成书于公元一世纪前后的《九章算术》是中国古代最重要的数学著作。其第八卷名为“方程”，
专门讨论线性方程组的解法。该卷的第一题是：

\begin{quote}
\kaishu
今有上禾三秉，中禾二秉，下禾一秉，实三十九斗；上禾二秉，中禾三秉，下禾一秉，实三十四斗；
上禾一秉，中禾二秉，下禾三秉，实二十六斗。问上、中、下禾实一秉各几何？
\end{quote}

这里“禾”指带秆的谷物，“秉”意为捆，“实”是打下的谷粒。设上、中、下三等禾每秉各出谷
$x$，$y$，$z$ 斗，题意就是
\begin{equation}\label{eq:he}
\left\{
\begin{aligned}
  3x+2y+\phantom{3}z&=39,\\
  2x+3y+\phantom{3}z&=34,\\
  x+2y+3z&=26.
\end{aligned}
\right.
\end{equation}

古人用算筹在算板上布列数字：每个方程占一列，自上而下依次是各未知数的系数和常数项，
三个方程自右向左排成三列（\cref{fig:suanchou}）。这样布列出来的数表呈方形，每一列称为
一“行”，这便是“方程”一名的由来。

\begin{figure}[htbp]
\centering
\begin{tikzpicture}[x=15mm, y=8.5mm, font=\kaishu]
  \fill[mjLight] (0,0) rectangle (3,4);
  \foreach \x in {1,2} \draw[mjRule, line width=.4pt] (\x,0) -- (\x,4);
  \foreach \y in {1,2,3} \draw[mjRule, line width=.4pt] (0,\y) -- (3,\y);
  \draw[mjInk, line width=.8pt] (0,0) rectangle (3,4);
  \foreach \c/\name in {0/左行, 1/中行, 2/右行}
    \node[font=\small\heitimd, mjAccent] at (\c+.5,4.45) {\name};
  \foreach \r/\name in {3/上禾, 2/中禾, 1/下禾, 0/实}
    \node[font=\small\heitimd, mjGray, anchor=west] at (3.12,\r+.5) {\name};
  \foreach \c/\a/\b/\cc/\d in {0/一/二/三/二十六, 1/二/三/一/三十四, 2/三/二/一/三十九} {
    \node at (\c+.5,3.5) {\a};
    \node at (\c+.5,2.5) {\b};
    \node at (\c+.5,1.5) {\cc};
    \node at (\c+.5,0.5) {\d};
  }
\end{tikzpicture}
\caption{“方程”第一题在算板上的布列：每列一个方程，自右向左排列}
\label{fig:suanchou}
\end{figure}

《九章算术》给出的解法称为\term[fang cheng shu]{方程术}，其要点是\term[bian cheng zhi chu]{遍乘直除}：
用右行上禾的系数 $3$ 遍乘中行各数，再从中行连续减去右行，直到中行顶端的数变为 $0$；对左行
作同样的处理。用今天的记号，把三个方程写成数表（每个方程占一行），这一步就是
\[
\left[\begin{array}{ccc|c}
3 & 2 & 1 & 39\\
2 & 3 & 1 & 34\\
1 & 2 & 3 & 26
\end{array}\right]
\xrightarrow[\;3r_3-r_1\;]{\;3r_2-2r_1\;}
\left[\begin{array}{ccc|c}
3 & 2 & 1 & 39\\
0 & 5 & 1 & 24\\
0 & 4 & 8 & 39
\end{array}\right]
\xrightarrow{\;5r_3-4r_2\;}
\left[\begin{array}{ccc|c}
3 & 2 & 1 & 39\\
0 & 5 & 1 & 24\\
0 & 0 & 36 & 99
\end{array}\right],
\]
其中 $3r_2-2r_1$ 表示“第二行乘以 $3$ 再减去第一行的 $2$ 倍”。最后一行即 $36z=99$，于是
$z=\frac{11}{4}$；代回第二行得 $y=\frac{17}{4}$，再代回第一行得 $x=\frac{37}{4}$。也就是说，
上禾一秉九斗四分斗之一，中禾一秉四斗四分斗之一，下禾一秉二斗四分斗之三——这与原书给出的
答案完全一致。

这个两千年前的算法，与今天计算机求解线性方程组所用的方法本质上是同一个。下面我们把它
整理成一般的理论。

\section{矩阵与初等行变换}\label{sec:matrix}

\begin{definition}\label{def:system}
含 $n$ 个未知数 $x_1,\dotsc,x_n$ 的 $m$ 个一次方程组成的方程组
\begin{equation}\label{eq:general-system}
\left\{
\begin{aligned}
  a_{11}x_1+a_{12}x_2+\dotsb+a_{1n}x_n&=b_1,\\
  a_{21}x_1+a_{22}x_2+\dotsb+a_{2n}x_n&=b_2,\\
  &\;\;\vdots\\
  a_{m1}x_1+a_{m2}x_2+\dotsb+a_{mn}x_n&=b_m
\end{aligned}
\right.
\end{equation}
称为\term[xian xing fang cheng zu]{线性方程组}。使每个方程都成立的数组 $(x_1,\dotsc,x_n)$ 称为它的一个
\term[jie]{解}，全体解组成的集合称为\term[jie ji]{解集}。有解的方程组称为\term[xiang rong]{相容}的。
若 $b_1=\dotsb=b_m=0$，则称方程组是\term[qi ci xian xing fang cheng zu]{齐次}的。
\end{definition}

方程组 \cref{eq:general-system} 的全部信息都包含在系数 $a_{ij}$ 与常数项 $b_i$ 中。把它们排成矩形数表，
\[
A=\begin{bmatrix}
a_{11} & a_{12} & \cdots & a_{1n}\\
a_{21} & a_{22} & \cdots & a_{2n}\\
\vdots & \vdots &        & \vdots\\
a_{m1} & a_{m2} & \cdots & a_{mn}
\end{bmatrix},
\qquad
(A\mid\vb{b})=\left[\begin{array}{cccc|c}
a_{11} & a_{12} & \cdots & a_{1n} & b_1\\
a_{21} & a_{22} & \cdots & a_{2n} & b_2\\
\vdots & \vdots &        & \vdots & \vdots\\
a_{m1} & a_{m2} & \cdots & a_{mn} & b_m
\end{array}\right],
\]
分别称为方程组的\term[xi shu ju zhen]{系数矩阵}与\term[zeng guang ju zhen]{增广矩阵}。一般地，$m$ 行 $n$ 列的
矩形数表称为 $m\times n$ \term[ju zhen]{矩阵}。

解方程组时我们对方程所做的操作，翻译成对增广矩阵的行所做的操作，就是下面三种变换。

\begin{definition}\label{def:row-op}
对矩阵施行的下列三种变换称为\term[chu deng hang bian huan]{初等行变换}：
\begin{enumerate}
  \item \textbf{对换}：交换两行的位置，记作 $r_i\leftrightarrow r_j$；
  \item \textbf{倍乘}：以非零数 $c$ 乘某一行的所有元素，记作 $cr_i$；
  \item \textbf{倍加}：把某一行的 $k$ 倍加到另一行上，记作 $r_j+kr_i$。
\end{enumerate}
若矩阵 $B$ 可以由 $A$ 经过有限次初等行变换得到，则称 $A$ 与 $B$ \term[hang deng jia]{行等价}，记作 $A\sim B$。
\end{definition}

每种初等行变换都可以用同类型的变换复原：$r_i\leftrightarrow r_j$ 的逆是它自己，$cr_i$ 的逆是
$\frac1cr_i$，$r_j+kr_i$ 的逆是 $r_j-kr_i$。因此行等价是一个等价关系。

\begin{theorem}\label{thm:row-equivalent}
若两个线性方程组的增广矩阵行等价，则这两个方程组同解。
\end{theorem}

\begin{proof}
只需对一次初等行变换证明。对换与倍乘显然不改变解集。对倍加变换 $r_j+kr_i$：若某数组满足
第 $i$、$j$ 两个方程，则它满足“第 $j$ 个方程加上第 $i$ 个方程的 $k$ 倍”；反之，新方程组的解满足
第 $i$ 个方程和新的第 $j$ 个方程，从而满足“新的第 $j$ 个方程减去第 $i$ 个方程的 $k$ 倍”，即原来的
第 $j$ 个方程。其余方程不变，故两个方程组同解。
\end{proof}

\section{阶梯形与高斯消元法}\label{sec:gauss}

消元的目标是把增广矩阵化为一种“阶梯”形状，从而可以自下而上逐个解出未知数。

\begin{definition}\label{def:echelon}
称矩阵是\term[hang jie ti xing]{行阶梯形}的，如果
\begin{enumerate}
  \item 全为零的行（如果有的话）都在矩阵的底部；
  \item 每个非零行的第一个非零元（称为该行的\term[zhu yuan]{主元}）都在上一行主元的右侧。
\end{enumerate}
若行阶梯形矩阵还满足：(3)~每个主元都是 $1$；(4)~主元所在列的其余元素都是 $0$，则称它是
\term[jian hua hang jie ti xing]{简化行阶梯形}的。
\end{definition}

下面两个矩阵分别是行阶梯形与简化行阶梯形的典型形状，其中 $\blacksquare$ 表示主元（非零），
$*$ 表示任意数：
\[
\begin{bmatrix}
\blacksquare & * & * & * & *\\
0 & \blacksquare & * & * & *\\
0 & 0 & 0 & \blacksquare & *\\
0 & 0 & 0 & 0 & 0
\end{bmatrix},
\qquad
\begin{bmatrix}
1 & 0 & * & 0 & *\\
0 & 1 & * & 0 & *\\
0 & 0 & 0 & 1 & *\\
0 & 0 & 0 & 0 & 0
\end{bmatrix}.
\]

\begin{theorem}\label{thm:rref}
任一矩阵都与唯一的一个简化行阶梯形矩阵行等价。
\end{theorem}

存在性由下面的算法给出；唯一性的证明留作习题（第~\ref{ex:rref-unique} 题）。

\begin{algorithm}[高斯--若尔当消元法]\label{alg:gauss-jordan}
\textbf{输入}：$m\times n$ 矩阵 $A$。\quad\textbf{输出}：与 $A$ 行等价的简化行阶梯形矩阵。
\begin{enumerate}[label=\arabic*., leftmargin=2em, itemsep=1pt]
  \item 若 $A$ 为零矩阵，转第 4 步。否则找出最左边的非零列，在其中取一个非零元作为主元，
        必要时对换两行，把主元所在的行换到第一行。
  \item 用倍加变换，把主元下方的元素全部化为零。
  \item 暂时“遮住”第一行，对余下的子矩阵重复第 1、2 步，直到矩阵成为行阶梯形。
        （以上称为\emph{前向消元}。）
  \item 自下而上，对每个主元：先用倍乘把它化为 $1$，再用倍加变换把它上方的元素化为零。
        （称为\emph{回代}。）
\end{enumerate}
\end{algorithm}

\begin{example}\label{ex:free-variables}
解线性方程组
\[
\left\{
\begin{aligned}
   x_1+2x_2-\phantom{3}x_3+3x_4&=4,\\
  2x_1+4x_2+\phantom{3}x_3+3x_4&=11,\\
  -x_1-2x_2+3x_3-5x_4&=-2.
\end{aligned}
\right.
\]
\end{example}

\begin{solution}
对增广矩阵施行高斯--若尔当消元：
\begin{align*}
&\left[\begin{array}{rrrr|r}
 1 &  2 & -1 &  3 &  4\\
 2 &  4 &  1 &  3 & 11\\
-1 & -2 &  3 & -5 & -2
\end{array}\right]
\xrightarrow[\;r_3+r_1\;]{\;r_2-2r_1\;}
\left[\begin{array}{rrrr|r}
 1 & 2 & -1 &  3 & 4\\
 0 & 0 &  3 & -3 & 3\\
 0 & 0 &  2 & -2 & 2
\end{array}\right]\\
&\xrightarrow[\;r_3-2r_2\;]{\;\frac13r_2\;}
\left[\begin{array}{rrrr|r}
 1 & 2 & -1 &  3 & 4\\
 0 & 0 &  1 & -1 & 1\\
 0 & 0 &  0 &  0 & 0
\end{array}\right]
\xrightarrow{\;r_1+r_2\;}
\left[\begin{array}{rrrr|r}
 1 & 2 & 0 &  2 & 5\\
 0 & 0 & 1 & -1 & 1\\
 0 & 0 & 0 &  0 & 0
\end{array}\right].
\end{align*}
主元位于第 1、3 列，对应的 $x_1$，$x_3$ 称为\term[zhu bian liang]{主变量}，其余的 $x_2$，$x_4$ 称为
\term[zi you bian liang]{自由变量}。令 $x_2=s$，$x_4=t$，得全部解
\[
\begin{bmatrix}x_1\\x_2\\x_3\\x_4\end{bmatrix}
=\begin{bmatrix}5\\0\\1\\0\end{bmatrix}
+s\begin{bmatrix}-2\\1\\0\\0\end{bmatrix}
+t\begin{bmatrix}-2\\0\\1\\1\end{bmatrix},
\qquad s,t\in\R .
\]
\end{solution}

\section{解的存在性与唯一性}\label{sec:existence}

\cref{ex:free-variables} 中的现象具有一般性。把增广矩阵化成简化行阶梯形以后，
方程组的解的情况可以一眼看出。

\begin{theorem}\label{thm:existence}
线性方程组相容，当且仅当其增广矩阵的行阶梯形中不出现形如
$\bigl[\,0\ \ 0\ \ \cdots\ \ 0\mid b\,\bigr]$（$b\ne0$）的行，即最后一列不是主元列。
在相容的情形下：
\begin{enumerate}
  \item 若没有自由变量，则方程组有唯一解；
  \item 若有自由变量，则方程组有无穷多个解。
\end{enumerate}
\end{theorem}

\begin{proof}
由\cref{thm:row-equivalent} 与\cref{thm:rref}，不妨设增广矩阵已是简化行阶梯形。若出现
$[\,0\ \cdots\ 0\mid b\,]$（$b\ne0$），它对应方程 $0=b$，方程组无解。否则，每个非零行都可以
把主变量用自由变量表示出来；给自由变量任意赋值，主变量便被唯一确定。于是方程组有解，
且解唯一当且仅当没有自由变量。
\end{proof}

\begin{corollary}\label{cor:homogeneous}
若齐次线性方程组中未知数的个数多于方程的个数，则它必有非零解。
\end{corollary}

\begin{proof}
齐次方程组总有零解，因而相容。主元个数不超过行数 $m$，而 $m<n$，所以至少有一个自由变量，
由\cref{thm:existence}，方程组有无穷多个解，其中必有非零解。
\end{proof}

\section{几何解释}\label{sec:geometry}

两个未知数的线性方程 $a_1x+a_2y=b$（$a_1$，$a_2$ 不全为零）在平面上表示一条直线，二元线性
方程组的解就是这些直线的公共点。两条直线的位置关系只有三种：相交、平行、重合，分别对应
唯一解、无解、无穷多解（\cref{fig:lines}）。三元线性方程组的情形类似，只是直线换成了空间中的平面。

\begin{figure}[htbp]
\centering
\begin{subfigure}{.31\textwidth}
\centering
\begin{tikzpicture}[scale=.68]
  \draw[mj axis] (-.4,0) -- (3.6,0); \draw[mj axis] (0,-.4) -- (0,3.4);
  \draw[mj curve] (-.2,3.1) -- (3.3,.3);
  \draw[mj curve2] (-.2,.4) -- (3.3,3.0);
  \node[mj point] at (1.62,1.74) {};
\end{tikzpicture}
\caption{相交：唯一解}
\end{subfigure}\hfill
\begin{subfigure}{.31\textwidth}
\centering
\begin{tikzpicture}[scale=.68]
  \draw[mj axis] (-.4,0) -- (3.6,0); \draw[mj axis] (0,-.4) -- (0,3.4);
  \draw[mj curve] (-.2,2.4) -- (3.3,.0);
  \draw[mj curve2] (-.2,3.3) -- (3.3,.9);
\end{tikzpicture}
\caption{平行：无解}
\end{subfigure}\hfill
\begin{subfigure}{.31\textwidth}
\centering
\begin{tikzpicture}[scale=.68]
  \draw[mj axis] (-.4,0) -- (3.6,0); \draw[mj axis] (0,-.4) -- (0,3.4);
  \draw[mj curve, line width=2.4pt, opacity=.35] (-.2,2.9) -- (3.3,.4);
  \draw[mj curve2] (-.2,2.9) -- (3.3,.4);
\end{tikzpicture}
\caption{重合：无穷多解}
\end{subfigure}
\caption{二元线性方程组的三种情形}
\label{fig:lines}
\end{figure}

\begin{remark}
\cref{thm:existence} 说明，线性方程组的解要么不存在，要么唯一，要么有无穷多个，不会恰好有
两个或三个。这一结论也可以直接看出：若 $\vb{u}\ne\vb{v}$ 都是解，则对一切实数 $t$，
$t\vb{u}+(1-t)\vb{v}$ 也是解（习题第~\ref{ex:two-solutions} 题）。
\end{remark}

\begin{exercises}
\exgroup{A 组}
\begin{enumerate}
  \item 用消元法解下列方程组：
  \par\smallskip
  \begin{tabular}{@{}l@{\qquad}l@{\qquad}l@{}}
    (1)~$\left\{\begin{aligned}x+y+z&=6\\x-y+2z&=5\\2x+y-z&=1\end{aligned}\right.$ &
    (2)~$\left\{\begin{aligned}x_1-x_2+2x_3&=1\\2x_1-x_2+3x_3&=4\\x_1+x_3&=3\end{aligned}\right.$ &
    (3)~$\left\{\begin{aligned}x+2y&=3\\2x+4y&=5\end{aligned}\right.$
  \end{tabular}
  \item 求矩阵 $\begin{bmatrix}1&2&3\\2&4&6\\1&0&1\end{bmatrix}$ 的简化行阶梯形。
  \item 讨论 $\lambda$ 取何值时，方程组
        $\left\{\begin{aligned}x+y+\lambda z&=1\\x+\lambda y+z&=1\\\lambda x+y+z&=1\end{aligned}\right.$
        有唯一解、无解、有无穷多解，并在有解时求出全部解。
  \item （百鸡问题）《张丘建算经》载：“今有鸡翁一，直钱五；鸡母一，直钱三；鸡雏三，直钱一。
        凡百钱买鸡百只，问鸡翁、母、雏各几何？”列出方程组，并求出它的全部非负整数解。
\end{enumerate}
\exgroup{B 组}
\begin{enumerate}
  \item 证明：两个线性方程组若同解，则它们的增广矩阵不一定行等价；但若两个\emph{相容}方程组
        同解且方程个数相同，它们的增广矩阵的简化行阶梯形相同。
  \item\label{ex:rref-unique} 证明\cref{thm:rref} 中的唯一性。（提示：对矩阵的列数归纳；
        若 $R$ 与 $S$ 是 $A$ 的两个简化行阶梯形，去掉最后一列后它们仍是同一矩阵的简化行阶梯形。）
  \item\label{ex:two-solutions} 证明：若线性方程组有两个不同的解，则它有无穷多个解。
  \item 设 $A\vb{x}=\vb{b}$ 相容，$\vb{x}_0$ 是它的一个解。证明：它的全部解恰好是
        $\vb{x}_0+\vb{y}$，其中 $\vb{y}$ 取遍齐次方程组 $A\vb{y}=\vb{0}$ 的全部解。
\end{enumerate}
\end{exercises}

\begin{history}
《九章算术》的“方程术”是世界上最早的系统的线性方程组解法\parencite{guo2009}。值得注意的是，
在消元过程中，系数难免出现负数，为此“方程”章同时提出了正负数的加减法则，即“正负术”：
“同名相除，异名相益，正无入负之，负无入正之。”这也是世界上关于负数运算的最早记载之一。
公元 263 年，刘徽为《九章算术》作注，对方程术的原理作了深入的阐释，还提出了减少运算
次数的改进方法。

在欧洲，类似的消元法直到 19 世纪初才被系统使用：高斯在用最小二乘法计算小行星轨道时，
需要求解大型的线性方程组，他所用的方法今天被称为高斯消元法。算法中“回代化简”的部分以
德国测地学家 W.~若尔当（Wilhelm Jordan，1842--1899）命名，他在测量学手册中推广了这一做法；
人们常把他与法国数学家 C.~若尔当混为一谈\parencite{althoen1987}。“矩阵”（matrix）一词则由
英国数学家西尔维斯特于 1850 年引入。
\end{history}
